\documentclass[11pt,reqno]{amsart}
\usepackage[margin=1.12in]{geometry}
\usepackage{microtype}
\usepackage{amssymb,amsmath,amsfonts,mathrsfs,mathtools}
\usepackage{xcolor}
\definecolor{mathblue}{HTML}{1F4E79}
\definecolor{mathteal}{HTML}{00656D}
\definecolor{mathpurple}{HTML}{643F7A}
\definecolor{mathgreen}{HTML}{2F6B3C}
\definecolor{mathurl}{HTML}{843C39}
\usepackage[colorlinks=true,linkcolor=mathblue,citecolor=mathgreen,urlcolor=mathurl,breaklinks=true]{hyperref}
\usepackage[nameinlink,capitalise,noabbrev]{cleveref}
\pdfmapfile{+cm.map}
\pdfmapfile{+cmextra.map}
\pdfmapfile{+symbols.map}
\pdfmapfile{+latxfont.map}
\pdfmapfile{+euler.map}
\pdfmapfile{+rsfs.map}
\newcommand{\Secref}[1]{{\hypersetup{linkcolor=mathteal}\Cref{#1}}}
\newcommand{\Eqref}[1]{{\hypersetup{linkcolor=mathpurple}\eqref{#1}}}
\newtheorem{theorem}{Theorem}[section]
\newtheorem{proposition}[theorem]{Proposition}
\newtheorem{lemma}[theorem]{Lemma}
\newtheorem{corollary}[theorem]{Corollary}
\theoremstyle{definition}
\newtheorem{definition}[theorem]{Definition}
\theoremstyle{remark}
\newtheorem{remark}[theorem]{Remark}
\crefname{theorem}{Theorem}{Theorems}
\crefname{lemma}{Lemma}{Lemmas}
\crefname{proposition}{Proposition}{Propositions}
\crefname{corollary}{Corollary}{Corollaries}
\crefname{definition}{Definition}{Definitions}
\newcommand{\Q}{\mathbf Q}
\newcommand{\K}{K}
\newcommand{\CH}{\operatorname{CH}}
\newcommand{\an}{\mathrm{an}}
\newcommand{\alg}{\mathrm{alg}}
\newcommand{\homol}{\mathrm{hom}}
\newcommand{\BC}{\mathrm{BC,act}}
\newcommand{\B}{\mathcal B}
\newcommand{\Obs}{\mathfrak o}
\newcommand{\id}{\operatorname{id}}
\newcommand{\Spec}{\operatorname{Spec}}
\newcommand{\Res}{\operatorname{Res}}
\newcommand{\pr}{\operatorname{pr}}
\newcommand{\Tr}{\operatorname{Tr}}
\newcommand{\cl}{\operatorname{cl}}
\newcommand{\Gal}{\operatorname{Gal}}
\newcommand{\Pic}{\operatorname{Pic}}
\newcommand{\codim}{\operatorname{codim}}
\newcommand{\et}{\mathrm{\acute{e}t}}
\newcommand{\M}{\mathscr M}
\numberwithin{equation}{section}
\title[Algebraic equivalence and strong Green forms]{Algebraic Equivalence and Chow--Motivic Transfer\\for Actual Non-Archimedean Green Forms}
\author{}
\date{Revised internal research draft, 16 November 2026}
\begin{document}
\begin{abstract}
Let $K/\Q_p$ be finite.  Given the Chow-invariant actual
Bott--Chern obstruction and the exact-support Green-pair construction
of the companion analytic theory, we prove that the obstruction
vanishes on geometrically algebraically trivial cycles of every
smooth projective $K$-variety.  The proof uses local rigidity of
actual forms in proper families and a compactness argument on
$p$-adic abelian groups.  Adding the numerical-model criterion and
correspondence covariance, we establish a transfer theorem for
homologically trivial cycles whose Chow classes factor through
abelian varieties after finite extension.  K\"unnemann's special-fiber
theorem supplies the required abelian input.  These deductions produce
zero-curvature Green forms on the precise complement of any
prescribed representative, tested against all compact actual
$\delta$-forms.  The four analytic inputs are stated as explicit
hypotheses so that the geometric deductions can be verified
independently of their companion proofs.
\end{abstract}
\maketitle
\pagestyle{plain}
\tableofcontents

\section{The two vanishing theorems}\label{sec:intro}

Let $K/\Q_p$ be a finite extension, with its fixed non-archimedean
absolute value, and let $X/K$ be smooth, projective and geometrically
integral of dimension $n$.  For a rational cycle $Z$ of codimension
$r$, $1\le r\le n$, put $U_Z=X^{\an}\setminus|Z|^{\an}$, where
$|Z|$ is the union of the components with nonzero coefficient.
Write $D=d'd''$ and denote by $\B^{p,q}$ the sheaf of ordinary actual
Mihatsch $\delta$-forms on the ordinary topology of $X^{\an}$
\cite{Mihatsch}.  The sought form is
\[
  \gamma_Z\in\B^{r-1,r-1}(U_Z),
  \qquad
  \int_{U_Z}\gamma_Z\wedge D\eta
          =-\int_{Z^{\an}}\eta
  \quad\bigl(\eta\in\B_c^{n-r,n-r}(X^{\an})\bigr).
\]
The integrals are those of actual $\delta$-forms; the construction
below provides their well-definedness.  In particular, the right
side retains the coefficients of the specified cycle.

The companion construction associates to every codimension-$r$ Chow
class a closed actual Bott--Chern class
\[
\Obs^r_{X,K}:\CH^r(X)_\Q\longrightarrow
 H^{r,r}_{\BC}(X):=
 \frac{\{\omega\in\B^{r,r}(X^{\an}):d'\omega=d''\omega=0\}}
 {D\B^{r-1,r-1}(X^{\an})}.
\]
Its zero class is a sufficient condition for the above Green form.
We state the precise companion results used in \Secref{sec:inputs}.

\begin{theorem}[Geometric algebraic equivalence]\label{thm:alg}
Assume the actual obstruction and exact-support properties
\textup{(A1)--(A2)} of \Cref{prop:inputs}.  For every
$z\in\CH^r(X)_{\Q,\alg}$, where algebraic equivalence is taken
over $\overline K$, one has
\[
                    \Obs^r_{X,K}(z)=0.
\]
For each specified rational cycle $Z$ representing $z$, there is a
form $\gamma_Z\in\B^{r-1,r-1}(U_Z)$ satisfying the preceding
zero-curvature identity for every compact actual test.
\end{theorem}

The proof in \Secref{sec:algproof} constructs an additive Chow family
over a Weil restriction of a Jacobian.  Its obstruction is an additive,
locally constant function on a compact $p$-adic abelian group and
hence vanishes.

We use rational correspondences with the convention that a degree-$a$
correspondence $V\rightsquigarrow W$ is a class in
$\CH^{\dim V+a}(V\times W)_\Q$.  Thus it sends codimension $r$
to codimension $r+a$.

\begin{definition}[Potential abelian Chow retraction]
\label[definition]{def:retraction}
The variety $X$ has a \emph{potential abelian Chow retraction} if
there are a finite extension $L/K$, abelian varieties $A_j/L$,
integers $a_j$, and correspondences
\begin{align*}
 u_j&\in\CH^{n+a_j}(X_L\times_L A_j)_\Q,\\
 v_j&\in\CH^{\dim A_j-a_j}(A_j\times_L X_L)_\Q
\end{align*}
for which
\begin{equation}\label{eq:diagonal}
        \Delta_{X_L}=\sum_j v_j\circ u_j
                   \quad\text{in }\CH^n(X_L\times_L X_L)_\Q.
\end{equation}
The zero-dimensional abelian variety $\Spec L$ is allowed.
\end{definition}

\begin{theorem}[Transfer through abelian Chow motives]
\label{thm:abelian}
Assume the companion properties \textup{(A1), (A2), (N), (C)} of
\Cref{prop:inputs}.  If $X$ has a potential abelian Chow retraction,
then for every $r$
\[
  \Obs^r_{X,K}(z)=0
         \qquad (z\in\CH^r(X)_{\Q,\homol}).
\]
For $1\le r\le n$, every specified rational cycle in such a class
has an ordinary actual strong Green form on its exact support
complement and against all compact actual tests.  For a fixed $r$,
it suffices that the right side of \Eqref{eq:diagonal} act as the
identity on $\CH^r(X_L)_{\Q,\homol}$.  The same conclusion holds
if its action is a fixed nonzero scalar multiple of the identity.
\end{theorem}

Here homological triviality means the vanishing of the geometric
$\ell$-adic cycle class for one (equivalently every) $\ell\ne p$.
The proof in \Secref{sec:motivproof} combines K\"unnemann's
semistable abelian model with the literal Chow diagonal.  This
separates the local analytic construction from the motivic
transfer, and keeps every operation on actual forms over the original
valued field.

Mihatsch constructs Green $\delta$-forms for complete-intersection
formal models \cite[Theorem~5.5]{Mihatsch}.  K\"unnemann established
unconditional height pairings for homologically trivial Chow cycles
on abelian varieties \cite{Kunnemann2001}.  The present statements
concern the actual differential-form representatives of local Green
identities.  Their analytic foundation is the internally developed
package described next.

\section{Actual curvatures, numerical models and correspondence kernels}
\label{sec:inputs}

This section records the exact companion results needed for the two
proofs.  They are internal analytic theorems, rather than assertions
from the external references cited for background.  The proofs
are contained in the frozen Rounds 10, 11, 12 and 16 \cite{R10,R11,R12,R16}.

\begin{proposition}[Companion analytic package]\label[proposition]{prop:inputs}
For smooth projective integral $K$-varieties $V$, the companion
construction has the following properties.

\smallskip\noindent
\textup{(A1)} There is a $\Q$-linear, Chow-invariant map
$\Obs^s_{V,K}:\CH^s(V)_\Q\to H^{s,s}_{\BC}(V)$, computed by
actual, separately closed degree-$s$ Chern-character forms of vector
bundles on projective flat models.  These model forms are functorial
under genuine analytic pullback, including rational fiber maps
$s_a:V\to V\times S$, with
$s_a^*\Obs^s_{V\times S,K}(\Theta)=\Obs^s_{V,K}(s_a^*\Theta)$.
If two such model characters have the same generic codimension-$s$
Chow class, their difference is $D$ of a global actual form.

\smallskip\noindent
\textup{(A2)} Every specified rational codimension-$s$ cycle $W$
has a pair
\[
 G_W\in\B^{s-1,s-1}(V^{\an}\setminus|W|^{\an}),\qquad
 \Omega_W\in\B^{s,s}(V^{\an})
\]
with $d'\Omega_W=d''\Omega_W=0$, $[\Omega_W]=\Obs^s_{V,K}([W])$,
and
\begin{equation}\label{eq:curved}
 \int_{V^{\an}\setminus|W|^{\an}}G_W\wedge D\eta
   +\int_{W^{\an}}\eta
    =\int_{V^{\an}}\Omega_W\wedge\eta
       \quad\bigl(\eta\in\B_c^{\dim V-s,\dim V-s}(V^{\an})\bigr).
\end{equation}

\smallskip\noindent
\textup{(N)} If $h:V'\to V$ is finite dominant, with $V'$ smooth
integral (possibly with a finite larger constant field), and a
projective flat model of $V'$ carries a rational virtual vector
bundle representing $h^*z$ whose degree-$s$ special-fiber Chern
character has degree zero on every integral vertical $s$-cycle,
then $\Obs^s_{V,K}(z)=0$.  Special-fiber multiplicities, residue
degrees and the full native degree of $h$ enter the criterion.

\smallskip\noindent
\textup{(C)} The kernels of the maps $\Obs^s_{V,K}$ are stable under
algebraic Chow correspondences between smooth projective integral
$K$-varieties, including proper algebraic pushforward.  For finite
$L/K$, the algebraic projection $\pi:V_L\to V$ satisfies
$\pi_*\pi^*=[L:K]$ on Chow groups.  The kernel assertion is
interpreted for $L$-schemes regarded as $K$-schemes, componentwise
after geometric base change.
\end{proposition}

The inputs split into two proof chains.  Properties (A1)--(A2)
come from model Chern forms, generic-Chow-zero primitives, and the
supported-perfect-complex Green-pair construction.  Property (N)
also uses the numerical detector for actual model Chern forms and
its finite-flat certificate descent.  Property (C) uses the actual
projective-bundle integration theorem and the smooth-immersion
deformation argument.  The latter is a statement about the kernel
of $\Obs$, not a form-valued pushforward along an arbitrary proper map.

\begin{lemma}[Correcting a prescribed Green pair]\label[lemma]{lem:correction}
Let $1\le s\le\dim V$.  Under \textup{(A2)}, if
$\Obs^s_{V,K}([W])=0$, then the zero-curvature Green form for $W$
can be chosen on precisely $V^{\an}\setminus|W|^{\an}$.
\end{lemma}
\begin{proof}
Choose $\Psi\in\B^{s-1,s-1}(V^{\an})$ with $D\Psi=\Omega_W$.
Then $G_W-\Psi|_{V^{\an}\setminus|W|^{\an}}$ satisfies
\Eqref{eq:curved} with zero right-hand side.  The correction is
global on $V^{\an}$, so the domain of the original $G_W$ is
unchanged.
\end{proof}

\section{A compact-parameter argument for algebraic equivalence}
\label{sec:compact}

The key geometric feature of actual forms is their dependence on
valuations of analytic units.  In a proper family, these valuations
remain literally constant when a rational parameter is sufficiently
close to a fixed point.

\begin{lemma}[Uniform constancy on rational fibers]
\label[lemma]{lem:rigid}
Let $S/K$ be smooth and separated, $a_0\in S(K)$ and
$\omega\in\B^{p,q}((X\times_K S)^{\an})$.  For
$s_a:X^{\an}\to(X\times S)^{\an}$ the section at $a\in S(K)$,
there is an ordinary analytic neighborhood $V$ of $a_0$ in $S^{\an}$
such that
\[
                     s_a^*\omega=s_{a_0}^*\omega
                 \qquad(a\in V\cap S(K)).
\]
\end{lemma}
\begin{proof}
Write $p:(X\times S)^{\an}\to X^{\an}$ and
$q:(X\times S)^{\an}\to S^{\an}$, and put $s=s_{a_0}$.
Locally near any point of $s(X^{\an})$, present $\omega$ as
$t_f^\star\kappa$ with $f=(f_1,\ldots,f_m)$ a tuple of analytic
units and $\kappa$ a Euclidean $\delta$-form.  On a common
neighborhood of the section, set $f_i^0=p^*s^*f_i$.  The difference
$f_i-f_i^0$ vanishes along the section.  Shrinking gives
\[
                      |f_i-f_i^0|<|f_i^0|,
\]
and hence $|f_i|=|f_i^0|$.  Thus $f$ and $f^0$ induce the same
tropical moment map on that neighborhood.  By the local definition
of actual forms and their functorial pullback
\cite[Definition~4.2 and Theorem~4.5]{Mihatsch},
$\omega=p^*s^*\omega$ there.

The resulting local neighborhoods cover $s(X^{\an})$; let $W$
denote their union.  Since $X$ is proper, $q$ is proper and closed.
The closed set $q((X\times S)^{\an}\setminus W)$ does not contain
$a_0$, because the analytic fiber over $a_0$ is exactly
$s(X^{\an})$.  Its complement is an ordinary neighborhood $V$
of $a_0$ with $q^{-1}(V)\subset W$.  Pull back
$\omega=p^*s^*\omega$ along each $s_a$ to obtain the assertion.
\end{proof}

\begin{lemma}[Multiples approach the identity]\label[lemma]{lem:multiple}
If $A/K$ is an abelian variety, $a\in A(K)$, and $U$ is a
neighborhood of $0$ in $A(K)$, then $ma\in U$ for some $m>0$.
\end{lemma}
\begin{proof}
$A(K)$ is compact, being a closed subset of $\mathbf P^N(K)$.
Choose a neighborhood $V$ of $0$ with $V-V\subset U$.
Finitely many translates of $V$ cover $A(K)$.  Two of sufficiently
many points $0,a,2a,\ldots$ lie in the same translate; their
difference is a positive multiple $ma\in U$.
\end{proof}

\begin{lemma}[Compact rigidity]\label[lemma]{lem:compact-rigidity}
Let $A$ be a compact topological abelian group and $W$ a torsion-free
abelian group.  An additive map $\varphi:A\to W$ is zero if its kernel
contains a neighborhood of the identity.
\end{lemma}
\begin{proof}
For any $a\in A$, the argument of \Cref{lem:multiple} gives a
positive integer $m$ such that $ma$ lies in the given neighborhood.
Then $m\varphi(a)=0$, and torsion-freeness gives
$\varphi(a)=0$.
\end{proof}

\begin{proposition}[Additive Chow families]\label[proposition]{prop:additive}
Assume \textup{(A1)}.  Let $A/K$ be an abelian variety and
$\Theta\in\CH^r(X\times A)_\Q$ satisfy
\[
 \Theta_0=0,\qquad
 \Theta_{a+b}=\Theta_a+\Theta_b
                      \quad(a,b\in A(K))
\]
for refined Chow fibers.  Then $\Obs^r_{X,K}(\Theta_a)=0$
for every $a\in A(K)$.
\end{proposition}
\begin{proof}
Represent $\Theta$ by the degree-$r$ Chern character of a rational
virtual bundle on $X\times A$, and extend the finitely many bundles
to a common projective flat model.  Let $\Omega$ denote its actual
model character.  Define
\[
       \varphi(a)=\Obs^r_{X,K}(\Theta_a)
                       =[s_a^*\Omega]\in H^{r,r}_{\BC}(X).
\]
The fiber identities and $\Q$-linearity of the obstruction make
$\varphi$ additive.  By \Cref{lem:rigid}, the actual forms
$s_a^*\Omega$ are constant for $a$ in a neighborhood of $0$.
Thus $\varphi$ is locally constant at the identity, and
$\varphi(0)=0$ since $\Theta_0=0$.  Its target is a real
vector space and hence torsion-free.  Apply
\Cref{lem:compact-rigidity} to the compact group $A(K)$.
\end{proof}

\section{The Poincar\'e family and proof of Theorem~\ref{thm:alg}}
\label{sec:algproof}

\begin{lemma}[Poincar\'e correspondence]\label[lemma]{lem:poincare}
Let $L/K$ be finite, let $C/L$ be a smooth projective geometrically
integral pointed curve, and let
$\Gamma\in\CH^r(X_L\times_L C)_\Q$.  With
$J=\Pic^0_{C/L}$, there is a class
$\Theta\in\CH^r(X_L\times_L J)_\Q$ such that
\[
 \Theta_t=\Gamma_*c_1(\mathcal P_t),\quad
 \Theta_0=0,\quad
 \Theta_{t+t'}=\Theta_t+\Theta_{t'}
                    \qquad(t,t'\in J(L)),
\]
where $\mathcal P$ is the normalized Poincar\'e line bundle.
\end{lemma}
\begin{proof}
On $X_L\times_L C\times_L J$ set
\[
 \Theta=(p_{XJ})_*
      \bigl(p_{XC}^*\Gamma\cdot p_{CJ}^*c_1(\mathcal P)\bigr).
\]
The projection has relative dimension one, so the resulting class
has codimension $r$.  Refined Gysin base change gives the asserted
fiber formula.  The normalized Poincar\'e bundle satisfies
$\mathcal P_{t+t'}\simeq\mathcal P_t\otimes\mathcal P_{t'}$,
and its zero fiber is trivial.  First Chern classes and Chow
correspondences are additive, proving the identities.
\end{proof}

\begin{lemma}[Norm of an additive family]\label[lemma]{lem:norm}
Let $L/K$ be finite Galois, $J/L$ an abelian variety, and
$\Theta\in\CH^r(X_L\times_L J)_\Q$ additive on $J(L)$.
Then $A=\Res_{L/K}J$ admits an additive family
$\Theta^N\in\CH^r(X\times_K A)_\Q$ whose fiber at
$a\in A(K)=J(L)$ corresponding to $t$ is
\[
                   (\Theta^N)_a=h_*\Theta_t,
                \quad h:X_L\longrightarrow X.
\]
\end{lemma}
\begin{proof}
Let $e:A_L\to J$ be evaluation and
$h_A:(X\times A)_L\to X\times A$ the finite flat projection.
Define
\[
 \Theta^N=(h_A)_*(\id_{X_L}\times e)^*\Theta.
\]
All operations here are algebraic Chow operations.  The
proper-pushforward/refined-Gysin base-change identity gives the
claimed formula for fibers.  Since $e$ is a group morphism, the
Chow fibers are additive.
\end{proof}

\begin{proof}[Proof of \Cref{thm:alg}]
Every geometrically algebraically trivial class is, after finite
extension, a finite rational sum of images of degree-zero divisors
on smooth projective curves under Chow correspondences.  This is
the curve-parameter description of algebraic equivalence
\cite{ACMV}.  Choose a finite Galois extension $L/K$ over which
the curves, correspondences, pointed structures, divisor classes,
and the rational-equivalence identities defining $z_L$ are all
defined.  Taking products of their Jacobians and sums of the
Poincar\'e families in \Cref{lem:poincare} gives an abelian variety
$J/L$, an additive family $\Theta$, and $t\in J(L)$ with
$\Theta_t=z_L$.

Apply \Cref{lem:norm}.  The resulting $K$-family is additive and
has a point $a\in A(K)$ such that
\begin{equation}\label{eq:norm}
           (\Theta^N)_a=h_*z_L=[L:K]z.
\end{equation}
By \Cref{prop:additive},
$[L:K]\Obs^r_{X,K}(z)=0$.  The target is a real vector
space, so division by $[L:K]$ proves $\Obs^r_{X,K}(z)=0$.
Finally, for the original specified representative $Z$, apply
\Cref{lem:correction}.  The global correction is performed on
$X^{\an}$, hence the complement remains $U_Z$.
\end{proof}

\section{Homologically trivial classes on abelian varieties}
\label{sec:abelianlocal}

We isolate the algebraic argument that converts K\"unnemann's
semistable model theorem into the numerical hypothesis of (N).
Let $F$ be a finite extension of $K$, $R=F^\circ$, and let
$\mathcal P/R$ be a regular projective strictly semistable model
of relative dimension $d$, with special fiber
$Y=\bigcup_i D_i$.  Write $Y^{(j)}$ for the disjoint union of the
$j$-fold component intersections.  Let
$A^q(Y)=\ker(A^q(Y^{(1)})\xrightarrow{\rho}A^q(Y^{(2)}))$
and
$A_q(Y)=\operatorname{coker}(A_{d-q-1}(Y^{(2)})
\xrightarrow{\gamma} A_{d-q}(Y^{(1)}))$, where the groups of
special-fiber cycles are taken modulo geometric homological
equivalence.

K\"unnemann proves for his semistable models of abelian varieties
that the geometric cycle class map on $A_q(Y)$ is injective and
that the complex
\begin{equation}\label{eq:kunn}
 A_{d+1-r}(Y)\xrightarrow{i^*i_*} A^r(Y)
               \xrightarrow{\cap[Y]} A_{d-r}(Y)
\end{equation}
is exact at the middle term
\cite[Theorems~5.6 and 6.6; \S5.8]{Kunnemann2001}.

\begin{lemma}[An extension with trivial component classes]
\label[lemma]{lem:extension}
Let $A/F$ be an abelian variety and let $\mathcal P/R$ be a
K\"unnemann semistable model with the preceding properties.
For $z\in\CH^r(A)_{\Q,\homol}$ there is
$\Theta_0\in\CH^r(\mathcal P)_\Q$ extending $z$ such that
its refined restrictions to all components $D_i$ have vanishing
geometric cycle classes.
\end{lemma}
\begin{proof}
Choose a cycle extension $\Theta$ and let $a\in A^r(Y)$ be
its tuple of component restrictions.  The class $a\cap[Y]$ is
the algebraic specialization of $z$.  Its $\ell$-adic class
vanishes because specialization commutes with the cycle class
map.  The injectivity on $A_{d-r}(Y)$ gives
$a\cap[Y]=0$.  Exactness of \Eqref{eq:kunn} supplies
$b\in A_{d+1-r}(Y)$ with $i^*i_*b=a$.
Lift $b$ to a rational Chow class $\beta$ on $Y$, and set
$\Theta_0=\Theta-i_*\beta$.  This has generic restriction $z$
and zero restriction in $A^r(Y)$.  Since $A^r(Y)$ is a subspace
of the direct sum of the component cycle-class images, each
component restriction has zero geometric class.
\end{proof}

\begin{proposition}[Abelian obstruction vanishing]
\label[proposition]{prop:abelianzero}
Assume \textup{(A1) and (N)} of \Cref{prop:inputs}.
For every abelian variety $A/F$, with $F/K$ finite, and every
$r\ge0$, the obstruction of every geometrically homologically
trivial class in $\CH^r(A)_\Q$ is zero when $A$ is regarded as
a smooth integral $K$-scheme.
\end{proposition}
\begin{proof}
The assertion is immediate in degree zero and impossible degrees.
Otherwise, enlarge $F$ finitely so that $A$ admits a regular
projective semistable model as above \cite{Kunnemann2001}.
By \Cref{lem:extension}, the pulled-back class admits an extension
whose restrictions to each geometric component of the special
fiber are homologically, hence numerically, trivial.

On the regular model, take the rational virtual-bundle class whose
Chern character represents this codimension-$r$ cycle and is
operationally pure.  Existence of this representative follows
from the rational Chern character isomorphism on regular schemes
\cite[Section~42.58]{StacksK}.  Its degree-$r$ operational
Chern class pairs to zero with every integral vertical $r$-cycle:
such a cycle is contained in a special-fiber component, and the
pairing is the numerical pairing with the refined restriction
there.  Regard the model over $K^\circ$.  Its scheme-theoretic
special-fiber multiplicities and residue-field degrees are
retained, and the finite morphism from the extended abelian
variety to the original has its full degree.  Property (N)
therefore gives zero obstruction in the original $K$-actual theory.
\end{proof}

\section{Proof of the Chow--motivic transfer theorem}
\label{sec:motivproof}

\begin{proposition}[Transfer from correspondence-generated classes]
\label[proposition]{prop:transfer}
Assume \textup{(A1) and (C)}.  Let $X/K$ be smooth projective
integral, and let $z\in\CH^r(X)_\Q$.  Suppose that after a finite
extension $L/K$ there are smooth projective $L$-varieties $V_j$,
classes $a_j\in\CH^{q_j}(V_j)_\Q$ of zero actual obstruction
over the original $K$, and rational correspondences
$v_j:V_j\rightsquigarrow X_L$ of degree $r-q_j$ such that
\[
              \lambda z_L=\sum_j(v_j)_*a_j
                      \quad\text{in }\CH^r(X_L)_\Q
\]
for some $\lambda\in\Q^\times$.  Then $\Obs^r_{X,K}(z)=0$.
\end{proposition}
\begin{proof}
Regard the $L$-varieties as smooth projective $K$-schemes.
By \textup{(C)}, the right side has zero obstruction on $X_L$.
For the finite projection $\pi:X_L\to X$, the algebraic identity
$\pi_*z_L=[L:K]z$ and \textup{(C)} give
$\lambda[L:K]z\in\ker\Obs^r_{X,K}$.
Since the kernel is a $\Q$-vector space, the assertion follows.
\end{proof}

\begin{proof}[Proof of \Cref{thm:abelian}]
Choose the data $L/K$, $u_j$, $v_j$ of
\Eqref{eq:diagonal}.  Fix $z\in\CH^r(X)_{\Q,\homol}$.
The classes
\[
               a_j:=(u_j)_*(z_L)
                  \in\CH^{r+a_j}(A_j)_\Q
\]
are geometrically homologically trivial, since geometric
cycle classes commute with correspondences.  By
\Cref{prop:abelianzero}, their native-$K$ obstruction classes
vanish.  The Chow diagonal identity supplies
$z_L=\sum_j(v_j)_*a_j$.  Apply \Cref{prop:transfer} with
$V_j=A_j$ and $\lambda=1$.  If the diagonal correspondence acts
as a nonzero scalar $\lambda$ on the indicated subgroup, the same
proposition applies with that scalar.
For $1\le r\le n$, the exact-support Green form follows from
\Cref{lem:correction}.  Only the action of
\Eqref{eq:diagonal} on the given Chow degree was used, proving
the fixed-degree assertion.
\end{proof}

\section{The algebraic--homological boundary}
\label{sec:boundary}

The first theorem is independent of the motivic hypothesis, whereas
the second acts on the larger Chow subgroup
\[
 \CH^r(X)_{\Q,\alg}\subseteq\CH^r(X)_{\Q,\homol}.
\]
The possible difference between the two groups is measured in part
by Griffiths-type phenomena.  On a variety with a potential abelian
Chow retraction, \Cref{thm:abelian} reaches such homologically
trivial classes as well.  Its proof uses two mechanisms beyond
the compact-parameter argument: numerical orthogonality on
K\"unnemann's special fibers, and preservation of the actual
obstruction kernel by the Chow correspondences in
\Eqref{eq:diagonal}.

Both vanishing theorems retain an explicit relationship to the
underlying cycle.  The model-character obstruction is computed
on a Chow class, but the final Green form is obtained from the
curved pair of the specified cycle.  Its singularity set is
therefore determined by that cycle, rather than by a convenient
representative chosen during the Chow calculation.

\begin{thebibliography}{99}
\bibitem{Mihatsch}
A.~Mihatsch, \emph{$\delta$-Forms on Lubin--Tate spaces},
arXiv:2112.10018, version 4 (2024), particularly Definition~4.2,
Theorems~4.5 and 5.5.

\bibitem{Kunnemann2001}
K.~K\"unnemann, \emph{Height pairings for algebraic cycles on
abelian varieties}, Ann. Sci. \`Ecole Norm. Sup. (4) \textbf{34}
(2001), 503--523, especially Theorems~5.6 and 6.6 and \S5.8.

\bibitem{ACMV}
J.~D.~Achter, S.~Casalaina-Martin and C.~Vial,
\emph{Parameter spaces for algebraic equivalence},
Int. Math. Res. Not. IMRN (2019), no.~6, 1863--1893.

\bibitem{StacksK}
The Stacks Project, \emph{Intersection Theory},
Section~42.58 (rational Chern character on regular schemes),
\url{https://stacks.math.columbia.edu/}.

\bibitem{R10}
Research archive, \emph{Model Chern forms without a finite trace},
internal source of Strong Green Round~10 (November 2026),
\texttt{chern/model\_chern\_without\_trace.tex}.

\bibitem{R11}
Research archive, \emph{The minimal finite trace and supported-complex
Green pairs}, internal sources of Round~11 (November 2026),
\texttt{one\_parameter\_trace.tex} and
\texttt{supported\_complex\_green.tex}.

\bibitem{R12}
Research archive, \emph{Compact parameters and numerical-model
descent}, internal sources of Round~12 (November 2026),
\texttt{compact\_parameter\_algebraic\_triviality.tex} and
\texttt{numerical\_witness\_descent.tex}.

\bibitem{R16}
Research archive, \emph{Immersion kernel covariance and proper
correspondence transfer}, internal sources of Round~16
(November 2026), including
\texttt{smooth\_immersion\_obstruction\_kernel.tex} and
\texttt{proper\_covariance\_and\_rational\_domination.tex}.
\end{thebibliography}
\end{document}
